\subsection{Two additional integral and spectral comparisons}

\begin{statement}{P1B.1}{Proved}
Let $N$ be an endomorphism of $\Z^6$ with $N^2=0$, rank three,
and primitive image. There is a unimodular basis in which
\[
 I+N=\begin{pmatrix}I_3&I_3\\0&I_3\end{pmatrix}.
\]
Its centralizer in $\operatorname{GL}_6(\Z)$ consists exactly of
$\left(\begin{smallmatrix}P&Q\\0&P\end{smallmatrix}\right)$,
with $P\in\operatorname{GL}_3(\Z)$ and $Q\in\operatorname{Mat}_3(\Z)$.
\end{statement}
\begin{proof}
The kernel of an integral homomorphism to a free group is
primitive. Both $\ker N$ and $\operatorname{im}N$ have rank three,
and the latter is contained in the former. Since the image is
primitive in $\Z^6$, their finite quotient is zero. Choose a basis
$e_1,e_2,e_3$ of this common group and choose $f_i$ with $Nf_i=e_i$.
The induced map $\Z^6/\ker N\to\operatorname{im}N$ is an
isomorphism, so the images of the $f_i$ form a quotient basis.
Thus $(e_1,e_2,e_3,f_1,f_2,f_3)$ is a unimodular basis and gives
the displayed matrix. This is also an algorithm: a Smith
decomposition gives the primitive kernel and solves the three
integral equations $Nf_i=e_i$.
Writing a potential centralizer element as
$\left(\begin{smallmatrix}P&Q\\R&S\end{smallmatrix}\right)$,
the equality $CN=NC$ gives $R=0$ and $S=P$.
Its determinant is $(\det P)^2$, so integral invertibility
is equivalent to $P\in\operatorname{GL}_3(\Z)$.
\end{proof}

\begin{statement}{P1B.2}{Proved}
A finite-order integral linear action on $\R^6$ admitting a
free finite affine torus lift has characteristic polynomial
among exactly thirty-eight cyclotomic products. Their linear
orders are
\[
 1,2,3,4,5,6,8,10,12.
\]
Exactly nineteen of these polynomials admit a commuting complex
structure on $\R^6$. These are a preliminary spectral list;
no assertion about global gluings is inferred from the list.
\end{statement}
\begin{proof}
The linear action has eigenvalue one. Otherwise $I-A$ is
invertible over $\R$, and the affine lift $x\mapsto Ax+b$
has a fixed point, obtained by solving $(I-A)x=b$ in $\R^6$
and projecting to the torus. A finite-order integral matrix is
semisimple and its characteristic polynomial is a product of
cyclotomic polynomials with nonnegative integer exponents.
The one-eigenvalue leaves at most five dimensions for the other
factors. A cyclotomic degree is even for every index greater
than two. The possible degrees are therefore one, two and four,
with indices respectively $1,2$; $3,4,6$; and $5,8,10,12$.
To verify that no index has been omitted, the formula
$\varphi(n)=\prod_{p^a\parallel n}p^{a-1}(p-1)$ with
$\varphi(n)\le4$ restricts $p$ to $2,3,5$, the powers to
$2^a$ with $a\le3$, $3^a$ with $a\le1$, and $5^a$ with
$a\le1$, and gives exactly those indices after multiplication.

The executed program \texttt{code/part1/spectral\_filter.py}
enumerates every exponent vector satisfying total degree six
and a positive exponent at $\Phi_1$, with bounds supplied by
these degrees. It writes all thirty-eight exponent vectors,
polynomials and least common multiples of their active indices.
Companion matrices in block diagonal form realize every listed
polynomial and order. A translation by a nonzero rational
fraction in a trivial eigenline makes a finite affine action
free: choose its denominator to be a multiple of the linear
order, and every nonidentity power has nonintegral translation
in that eigenline.

A commuting real complex structure requires even multiplicity
of each real eigenvalue $1$ and $-1$. Conversely, on these real
eigenspaces pair basis vectors into oriented complex lines;
on every conjugate nonreal eigenpair use its rotation complex
structure. This constructs a commuting complex structure.
The program checks these two parity conditions and obtains
nineteen products. Since the remaining cyclotomic factors have
even degrees and the total degree is six, the two real
multiplicities have the same parity. Thus the nineteen
complex-compatible products are exactly the orientation-preserving
ones in this preliminary list.
\end{proof}
